\begin{frame}{역사: 부스터(1990) $\to$ AdaBoost(1997) $\to$ XGBoost(2016)}
  \small
  \begin{tabular}{@{}llp{8.2em}@{}}
    \toprule
    연도 & 이름 & 내용 \\
    \midrule
    1990 & booster & Schapire가 ``순차 추가''의 \textbf{부스팅}이라는
      이름 붙임 \\
    1997 & AdaBoost & Freund \& Schapire: \textbf{틀린 샘플의 가중치
      올려} 다음 모델이 그쪽에 집중 $\to$ 첫 실용 알고리즘 \\
    2001 & GBDT & Friedman: ``부스팅 = \textbf{함수 공간의 경사하강}''
      해석 $\to$ 명확한 이론적 지위, 어떤 손실에든 적용할 일반 공식 \\
    2016 & XGBoost & Chen \& Guestrin: \textbf{이차 타원 전개}(2차
      미분 근사) + 정규화 항 + 분기 탐색 최적화 $\to$ 2010년대 중후반
      Kaggle 정형 데이터 쓸어담음 \\
    2017 & LightGBM & ``위에서 아래로 자란다'' 성장 전략 + 특징 압축으로
      \textbf{속도 극대화} (Microsoft) \\
    \bottomrule
  \end{tabular}
  \vskip1em
  {\footnotesize ``정형 데이터라면 \textbf{먼저 GBDT를 돌려보라}''가
  업계의 기본 반사가 된 순간. scikit-learn \texttt{GradientBoosting\_*} =
  Friedman 원형, XGBoost/LightGBM/CatBoost = ``더 빠르고 튜닝 편한''
  세 주요 변주.}
\end{frame}
